% Plain TeX document
% Review of APS_repair.pdf against ch24_revisions_claude.pdf
% Compile with: tex aps_repair_review_plain.tex

\input texsis
\input paper
%\magnification=\magstep1
%\baselineskip=14pt
%\parskip=6pt
%\parindent=20pt
%\hsize=5.5in
%\vsize=8.5in
%\hoffset=0.75in
%
%\font\titlefont=cmr12 scaled\magstep2
%\font\sectionfont=cmbx12
%\font\subsecfont=cmbx10
%\font\smalrm=cmr9
%
%\def\thedate{\number\day\space
%  \ifcase\month\or January\or February\or March\or April\or May\or June\or
%  July\or August\or September\or October\or November\or December\fi
%  \space\number\year}
%
%\def\bigtitle#1{\vskip12pt\centerline{{\titlefont #1}}\vskip4pt}
%\def\section#1{\vskip14pt\noindent{\sectionfont #1}\vskip4pt\noindent}
%\def\subsec#1{\vskip10pt\noindent{\subsecfont #1}\vskip4pt\noindent}
%\def\bquote{\vskip4pt\leftskip=20pt\rightskip=10pt}
%\def\ebquote{\vskip4pt\leftskip=0pt\rightskip=0pt}
%\def\bitem{\vskip2pt\noindent$\bullet$\enspace}
%
%\bigtitle{Review of {\tt APS\_repair.pdf}:}
%\bigtitle{Remaining Errors and Recommendations}
%\centerline{Review against {\tt ch24\_revisions\_claude.pdf}}
%\vskip4pt\centerline{\thedate}\vskip14pt
%
\noindent{\bf Abstract.}
The document {\tt APS\_repair.pdf} successfully incorporates the principal
logical corrections recommended in {\tt ch24\_revisions\_claude.pdf}:
Lemma~2, Definitions~7 and~8, the monotonicity proof, Theorem~1, the
compactness proposition, and the two-value corollary are all correctly
stated and correctly ordered.  Three expository residues of the old
four-tuple formulation remain, and one algorithmic passage should either
be rewritten or given an explicit citation to the corollary that licenses
it.  This note identifies each issue and states the required correction.

\section{1. Background}

The core repair identified in {\tt ch24\_revisions\_claude.pdf} is the
replacement of every admissible {\it four-tuple\/} $(x,y,w_1,w_2)$ by an
admissible {\it triple\/} $(x,y,w(\cdot))$, where $w(\cdot):Y\to W$ is a
function, throughout the argument that establishes compactness of $V$.
The two-value characterization---using scalar continuation values
$(v_1,v_2)\in V\times V$---is then recovered as a {\it corollary\/} once
compactness is in hand.  The logical chain is:
$$\hbox{Lemma 2 (triples)}\;\Rightarrow\;
  \hbox{Defs.\ 7--8 (triples)}\;\Rightarrow\;
  \hbox{Theorem 1}\;\Rightarrow\;
  \hbox{Prop.\ 2 (compactness)}\;\Rightarrow\;
  \hbox{Corollary (two-value).}$$

\section{2. Errors That Remain in {\tt APS\_repair.pdf}}

\subsec{2.1. Four-tuple language in the paragraph following Definition 7
(p.~1141)}

After stating Definition~7 correctly in triple form, the immediately
following paragraph reverts to the old language:

\bquote
``Notice that when $W\subset V$, the admissible {\bf 4-tuple}
$(x,y,w_1,w_2)$ determines an SPE with strategy profile
$\sigma_1=(x,y)$, $\sigma|_{(x,y)}=\sigma^1$,
$\sigma|_{(\nu,\eta)}=\sigma^2$ for $(\nu,\eta)\ne(\tilde x,\tilde y)$,
where $\sigma^1$ is a continuation strategy that yields value
$w_1=V_g(\sigma^1)$ and $\sigma^2$ yields continuation value
$w_2=V_g(\sigma^2)$.  The value of the SPE is
$V_g(\sigma)=w=(1-\delta)r(x,y)+\delta w_1$.''
\ebquote

This paragraph introduces scalar continuation values $w_1$ and $w_2$
under the label ``admissible 4-tuple,'' contradicting the corrected
Definition~7 that immediately precedes it.

\vskip6pt\noindent{\it Recommendation 1.}\enspace
Replace the offending paragraph with the following:

\bquote
``Notice that when $W\subset V$, an admissible triple $(x,y,w(\cdot))$
determines an SPE with strategy profile
$$\sigma_1=(x,y),\qquad \sigma|_{(x,\eta)}=\tilde\sigma(\eta)
  \quad\forall\,\eta\in Y,$$
where $\tilde\sigma(\eta)$ is any SPE with value $w(\eta)\in V$.
The value of the SPE is
$V_g(\sigma)=(1-\delta)\,r(x,y)+\delta\,w(y)$.''
\ebquote

\subsec{2.2. Two occurrences of ``4-tuples'' in Section 25.8.1
(p.~1141--1142)}

Section~25.8.1 (``Basic idea of dynamic programming squared'') contains
two sentences that still use four-tuple language:

\vskip4pt
\noindent (a) ``APS define a mapping $B(W)$ that, by considering only
admissible {\bf 4-tuples}, maps a {\it set\/} of values $W$ tomorrow
into a new {\it set\/} $B(W)$ of values today.''

\vskip2pt
\noindent (b) ``Thus, APS use admissible {\bf 4-tuples} to map candidate
continuation values tomorrow into new candidate values today.''

\vskip6pt\noindent{\it Recommendation 2.}\enspace
Replace both occurrences of ``4-tuples'' with ``triples'':

\vskip4pt
\noindent (a) ``APS define a mapping $B(W)$ that, by considering only
admissible {\bf triples}, maps a set of values $W$ tomorrow into a new
set $B(W)$ of values today.''

\vskip2pt
\noindent (b) ``Thus, APS use admissible {\bf triples} to map candidate
continuation values tomorrow into new candidate values today.''

\subsec{2.3. The iteration algorithm in Section 25.9 uses the two-value
form without citing the corollary (p.~1145--1146)}

The algorithm for computing the boundaries of $B(W_0)$ states the
problem for the worst value $\underline w_1$ as a minimization over
$(x,y)\in C$ and $(w_1,w_2)\in[\underline w_0,\overline w_0]^2$:
$$\underline w_1
  = \min_{\scriptstyle(x,y)\in C;\;(w_1,w_2)\in[\underline w_0,
    \overline w_0]^2}
    (1-\delta)\,r(x,y)+\delta w_1,$$
subject to
$$(1-\delta)\,r(x,y)+\delta w_1\;\geq\;
  (1-\delta)\,r(x,\eta)+\delta w_2\quad\forall\,\eta\in Y.$$

This is the four-tuple formulation.  It is mathematically correct
{\it given\/} the two-value Corollary, but in the revised chapter the
Corollary appears {\it before\/} this algorithm (p.~1144--1145).  The
algorithm must therefore include an explicit reference to the Corollary
that licenses the reduction from a function $w(\cdot)$ to two scalar
values $(w_1,w_2)$.  As currently written it does neither.

\vskip6pt\noindent{\it Recommendation 3.}\enspace
Add a bridging sentence immediately before the two optimization problems:

\bquote
``By the two-value Corollary above, it suffices to search over pairs
$(w_1,w_2)\in W\times W$ rather than full functions $w(\cdot):Y\to W$;
the boundary problems therefore reduce to:''
\ebquote

\noindent Alternatively, rewrite the objective in triple form
$$\underline w_1
  = \min_{\scriptstyle(x,y,w(\cdot))\;\hbox{\smalrm admissible w.r.t.}\;
    [\underline w_0,\overline w_0]}
    (1-\delta)\,r(x,y)+\delta\,w(y),$$
and note that by the Corollary the minimum is attained at a $w(\cdot)$
taking at most two values.

\section{3. Summary of Remaining Issues}

\vskip6pt
\noindent
\halign{%
  \vrule\quad\strut#\hfil\quad&\vrule\quad#\hfil\quad&
  \vrule\quad#\hfil\quad\vrule\cr
\noalign{\hrule}
{\bf Location}&{\bf Issue}&{\bf Required action}\cr
\noalign{\hrule}
p.~1141, after Def.~7&
  ``4-tuple $(x,y,w_1,w_2)$''&
  Replace with triple description\cr
\noalign{\hrule}
p.~1141--42, Section~25.8.1 ($\times 2$)&
  ``admissible 4-tuples'' (two places)&
  Replace both with ``admissible triples''\cr
\noalign{\hrule}
p.~1145--46, algorithm in Section~25.9&
  $(w_1,w_2)\in W^2$ without justification&
  Add citation to two-value Corollary\cr
\noalign{\hrule}}

\section{4. Items Correctly Implemented}

The following elements of {\tt ch24\_revisions\_claude.pdf} are correctly
incorporated in {\tt APS\_repair.pdf}:

\bitem {\bf Lemma~2} (p.~1140): states and proves the necessary and
sufficient condition using a function $v(\cdot):Y\to V$.

\bitem {\bf Definition~7} (p.~1141): correctly defines admissible
{\it triples\/} $(x,y,w(\cdot))$.

\bitem {\bf Definition~8} (p.~1142): $B(W)$ defined via admissible triples.

\bitem {\bf Monotonicity proof} (p.~1142--43): uses the triple/function
formulation throughout.

\bitem {\bf Theorem~1 proof} (p.~1143--44): constructs sequences of
functions $w_1(\cdot),w_2(\cdot),\ldots$ correctly.

\bitem {\bf Proposition~2} (compactness, p.~1144): present, correctly
placed, and correctly proved.

\bitem {\bf Corollary} (two-value characterization, p.~1144--45):
correctly positioned {\it after\/} compactness, with a correct necessity
proof using $v_2=\inf_{\eta\in Y}v(\eta)\in V$.

\vfill\eject
\bye
